%! Author = denis
%! Date = 20/08/2026

\begin{frame}{Circuitos Combinacionais: Estrutura Geral}

    Um circuito combinacional pode ser representado por:

    \[
        \boxed{\text{Entradas}}
        \longrightarrow
        \boxed{\text{Circuito Lógico}}
        \longrightarrow
        \boxed{\text{Saídas}}
    \]

    \vspace{0.5cm}

    Exemplo:

    \[
        A,B,C
        \longrightarrow
        \boxed{\text{AND, OR, NOT}}
        \longrightarrow
        F
    \]

    \vspace{0.5cm}

    \begin{alertblock}{Importante}
        O circuito não possui memória.

        Se as entradas mudarem, a saída também poderá mudar imediatamente.
    \end{alertblock}

\end{frame}